% Part VI: where the construction stands in the field.
% ASCII source only. \input after part5.tex.
\newpage
\proofpart{6}
\pagetitle{26 / WHERE THIS STANDS IN THE FIELD}{What already rhymes with the consensus, and what does not}

Most of my construction is not new. Its lattice layer is Wilson's, its two-dimensional layer is exact Yang--Mills as
Migdal, Gross--Witten and Driver--Sengupta--L\'evy left it, its small room is the finite-volume programme of L\"uscher and
van Baal moved to a stretched sphere, and its renormalization is one-loop asymptotic freedom. What is mine is the
organisation by bundle layers, the Helike carrier, the sources that hold the background, and the claim that the lowest
note is inherited along scale. The table says which is which, so that nothing standard is presented as new and nothing new
is presented as standard.

\begin{center}\small\setlength{\tabcolsep}{4pt}\renewcommand{\arraystretch}{0.97}
\begin{tabular}{@{}p{2.2in}p{3.5in}p{1.1in}@{}}\toprule
\textbf{Part / result} & \textbf{Established line it continues} & \textbf{Standing}\\\midrule
0. Main Conjecture (i) & Jaffe--Witten Clay statement; Osterwalder--Schrader axioms & standard\\
0. Main Conjecture (ii), the inherited note & strong--weak crossover of 4-D $SU(2)$ lattice theory (Creutz 1980) & departure\\
0. Bundle-layer organisation & standard principal-bundle geometry of gauge fields & reformulation\\\midrule
I. Wilson loops, Giles, measure on $\mathcal A/\mathcal G$ & Wilson 1974; Giles 1981; Gribov 1978; Singer 1978 & standard\\
I. Reflection positivity, strong-coupling gap & Osterwalder--Seiler 1978; Seiler 1982 & standard\\
I. Electric gap $3g^2/(2a)$ on one plaquette & Kogut--Susskind 1975 strong-coupling Hamiltonian & standard\\
I. ``Leading term of the full gap'' & strong-coupling expansions (Drouffe--Zuber 1983 review) & reformulation\\\midrule
II. 2-D area law, every spin; Casimir scaling & Migdal 1975; Gross--Witten 1980; Driver 1989; Sengupta; L\'evy & standard\\
II. Centre vortex flips half-integer loops & 't Hooft 1978; Greensite 2003 review & standard\\
II. GIGI checks & planted exact answers & new but consistent\\\bottomrule
\end{tabular}
\end{center}

\begin{center}\small\setlength{\tabcolsep}{4pt}\renewcommand{\arraystretch}{0.97}
\begin{tabular}{@{}p{2.2in}p{3.5in}p{1.1in}@{}}\toprule
\textbf{Part / result} & \textbf{Established line it continues} & \textbf{Standing}\\\midrule
III. $\mathbb R^9/SO(3)$, vacua, $CS(\lambda)$, sphaleron & modes on $S^3$ (van Baal--van den Heuvel 1994); BPST 1975; Manton 1983 & standard\\
III. Theorem A, discrete spectrum & Schr\"odinger operators (Reed--Simon) & standard\\
III. Helike as Hessian floor (\emph{By Enough}) & Hessians about homogeneous fields on $S^3$ & new but consistent\\
III. Sourced reduced model & none: sources holding a non-vacuum background & departure\\
III. $S_C$, the ordering & computer-assisted eigenvalue bounds (Kato--Temple, Lehmann) & new but consistent\\
III. Fold conjecture & Agmon 1982; Helffer--Sj\"ostrand 1984 & new but consistent\\\midrule
IV. Fast fibre, flat Berry phase, empty sphere & Born--Oppenheimer on $S^3$ (van den Heuvel 1996); Ikeda--Taniguchi 1978 & standard\\
IV. $\lambda=0\sim1$; cutoff must be spectral; shift $-2$ & large gauge transformations; gauge-invariant regularization & new but consistent\\
IV. Log divergence $=b_0/8\pi^2$ & Gross--Wilczek, Politzer 1973; heat-kernel $a_4$ (Vassilevich 2003) & reformulation\\
IV. M3, uniform dressing, $g$ adiabatic & L\"uscher 1983; Koller--van Baal 1988 (torus); Ba{\l}aban & departure\\\midrule
V. Exact crossing across charge sectors & Kato--Rellich analytic perturbation theory & standard\\
V. Switch $s^*$, carrier persistence (M8) & level tracking in finite-volume spectra & new but consistent\\
V. Glueball link & lattice $0^{++}<2^{++}$ (Morningstar--Peardon 1999) & departure\\
V. Agmon length; $\theta$ Bloch phase; $e^{-8\pi^2/g^2}$ & 't Hooft 1976; Jackiw--Rebbi; Callan--Dashen--Gross 1976 & standard\\
V. Source bridge; M2, the infrared gap & none; M2 is the Clay problem itself & departure\\
V. UV results; RP survives limits & Osterwalder--Seiler; Ba{\l}aban; Magnen--Rivasseau--S\'en\'eor 1993 & standard\\\bottomrule
\end{tabular}
\end{center}

\textbf{Where a referee will push.} (1) The \emph{sources}: Helike is lowest only because the sources hold the background
past the vacuum; without them a charge-$0$ singlet is lowest (Part III). Nothing in the literature uses such a sourced
reduction, and the source bridge that would remove them is open. (2) The \emph{slow--fast split}: on the round sphere
the lowest modes at $A=0$ are eighteen, not nine; van Baal and van den Heuvel keep both the left- and right-invariant
families because they are degenerate. My split puts the second family in the fast fibre and rests the separation on $g$,
not on a frequency gap, which is unconventional and is exactly what M3 must justify. (3) The \emph{glueball link}: a
squashed sphere at $t=3$ breaks rotations, so continuing charge labels to $J^{PC}$ in infinite volume needs an argument
I do not have. (4) The \emph{inherited note}: lattice lore supports a smooth crossover from strong to weak coupling, but
nobody has traced a single carrier through it, and the carrier changes identity even in my own small room. (5) M2 is the
infrared half of the Clay problem; I claim no progress on it. (6) Helike rests on my own Zenodo papers, not yet refereed.

So the proved parts are standard or reformulated, the new parts are labelled as conjecture or computation, and the
departures are the ones listed above, named before anyone has to find them.
